\setchapterpreamble[u]{\margintoc}
\chapter[Пути с ограничениями в терминах формальных языков]{Задача о поиске путей с ограничениями в терминах формальных языков}
\label{chpt:FLPQ}
\tikzsetfigurename{FLPQ_}

В данной главе сформулируем постановку задачи о поиске путей в графе с ограничениями.
Также мы приведём несколько примеров областей, в которых применяются алгоритмы решения этой задачи.

\section{Постановка задачи}

Пусть нам дан конечный ориентированный помеченный граф $\mscrG = \langle V, E, L \rangle$.
Функция $\omega(\pi) = \omega((v_0, l_0, v_1), (v_1, l_1, v_2), \dots, (v_{n-1}, l_{n-1}, v_n)) = l_0 \cdot l_1 \cdot \dots \cdot l_{n-1}$ строит слово по пути посредством конкатенации меток рёбер вдоль этого пути.
Очевидно, для пустого пути данная функция будет возвращать пустое слово, а для пути длины $n  > 0$~--- непустое слово длины $n$.

Если теперь рассматривать задачу поиска путей, то окажется, что то множество путей, которое мы хотим найти, задаёт множество слов, то есть язык.
А значит, критерий поиска мы можем сформулировать следующим образом: нас интересуют такие пути, что слова, составленные из меток вдоль них, принадлежат заданному языку.

\begin{definition}[Задача поиска путей с ограничениями в терминах формальных языков]
    \label{def1}
    \emph{Задача поиска путей с ограничениями в терминах формальных языков} заключается в поиске множества путей $\Pi = \{\pi \mid \omega(\pi) \in \mscrL\}$.
\end{definition}

В задаче поиска путей мы можем накладывать дополнительные ограничения на путь (например, чтобы он был простым, кратчайшим или Эйлеровым~\sidecite{kupferman2016eulerian}), но это уже другая история.

Другим вариантом постановки задачи является задача достижимости.

\begin{definition}[Задача достижимости]
    \label{def2}
    \emph{Задача достижимости} заключается в поиске множества пар вершин, для которых найдется путь с началом и концом в этих вершинах, что слово, составленное из меток рёбер пути, будет принадлежать заданному языку
    \[\Pi' = \{(v_{i}, v_{j}) \mid \exists v_{i} \pi v_{j}, \omega(\pi) \in \mscrL\}.\]
\end{definition}

При этом, множество $\Pi$ может являться бесконечным, тогда как $\Pi'$ конечно, по причине конечности графа $\mscrG$.

Язык $\mscrL$ может принадлежать разным классам и быть задан разными способами.
Например, он может быть регулярным, контекстно-свободным, или многокомпонентным контекстно-свободным.

Если $\mscrL$~--- регулярный, $\mscrG$ можно рассматривать как недетерминированный конечный автомат (НКА), в котором все вершины являются одновременно и стартовыми, и конечными.
Тогда задача поиска путей, в которой $\mscrL$~--- регулярный, сводится к пересечению двух регулярных языков.

Более подробно мы рассмотрим случай, когда $\mscrL$~--- контекстно-свободный язык.

Путь $G = \langle \Sigma, N, P \rangle$~--- контекстно-свободная грамматика.
Будем считать, что $L \subseteq \Sigma$.
Мы не фиксируем стартовый нетерминал в определении грамматики, поэтому, чтобы описать язык, задаваемый ей, нам необходимо отдельно зафиксировать стартовый нетерминал.
Таким образом, будем говорить, что $L(G,N_i) = \{ w \mid N_i \xRightarrow[G]{*} w  \}$~--- это язык задаваемый грамматикой $G$ со стартовым нетерминалом $N_i$.

\begin{example}
    Пример задачи поиска путей.

    Дана грамматика $G$, задающая язык $\mscrL = a^n b^n$:
    \begin{align*}
        S & \to a b   \\
        S & \to a S b
    \end{align*}
    И дан граф $\mscrG$, представленный на рисунке~\ref{fig:all_paths_example}.
\begin{marginfigure}
    \begin{center}
        \input{figures/edge_lebelled_graph_a_3_loop_b_2_loop}
    \end{center}
    \caption{Входной граф}
    \label{fig:all_paths_example}
\end{marginfigure}

    Кратчайшими путями, принадлежащими множеству $\Pi = \{\pi \mid \omega(\pi) \in \mscrL\}$, являются, например, пути, представленные на изображении~\ref{fig:paths}.

\begin{marginfigure}
    \begin{center}
        \input{figures/flpq/path1.tex}
        \input{figures/flpq/path2.tex}
    \end{center}
    \caption{Примеры путей}
    \label{fig:paths}
\end{marginfigure}

\end{example}


\section{О разрешимости задачи}

Задачи из определения~\ref{def1} и~\ref{def2} сводятся к построению пересечения языка $\mscrL$ и языка, задаваемого путями графа, $R$.
А мы для обсуждения разрешимости задачи рассмотрим более слабую постановку задачи:

\begin{definition}[TODO: Что ты?]
    Необходимо проверить, что существует хотя бы один такой путь $\pi$ для данного графа, для данного языка $\mscrL$, что $\omega(\pi) \in \mscrL$.
\end{definition}

Эта задача сводится к проверке пустоты пересечения языка $\mscrL$ c $R$~--- регулярным языком, задаваемым графом.
От класса языка $\mscrL$ зависит её разрешимость:
\begin{itemize}
    \item Если $\mscrL$ регулярный, то получаем задачу пересечения двух регулярных языков: $\mscrL \cap R = R'$.
          $R'$~--- также регулярный язык.
          Проверка регулярного языка на пустоту~--- разрешимая проблема.
    \item Если $\mscrL$ контекстно-свободный, то получаем задачу: $\mscrL \cap R = CF$~--- контекстно-свободный.
          Проверка контекстно-свободного языка на пустоту~--- разрешимая проблема.
    \item Помимо иерархии Хомского существуют и другие классификации языков.
          Так например, класс конъюнктивных (Conj) языков~\sidecite{DBLP:journals/jalc/Okhotin01}является строгим расширением контекстно-свободных языков и все так же позволяет полиномиальный синтаксический анализ.

          Пусть $\mscrL$~--- конъюнктивный.
          При пересечении конъюнктивного и регулярного языков получается конъюнктивный ($\mscrL \cap R = Conj$), а проблема проверки Conj на пустоту не разрешима~\sidecite{DBLP:journals/tcs/Okhotin03a}.
    \item Ещё один класс языков из альтернативной иерархии, не сравнимой с Иерархией Хомского,~--- MCFG (multiple context-free grammars)~\sidecite{SEKI1991191}.
          Как его частный случай~--- TAG (tree adjoining grammar)~\sidecite{Joshi1997}.

          Если $\mscrL$ принадлежит классу MCFG, то $\mscrL \cap R$ также принадлежит MCFG.
          Проблема проверки пустоты MCFG разрешима~\sidecite{SEKI1991191}.
\end{itemize}

Существует ещё много других классификаций языков, но поиск универсальной иерархии до сих пор продолжается.

Далее, для изучения алгоритмов решения, нас будет интересовать задача $R \cap CF$.

\section{Области применения}

Поиск путей с ограничениями в виде формальных языков широко применяется в различных областях.
Ниже даны ключевые работы по применению поиска путей с контекстно-свободными ограничениями и ссылки на них для более детального ознакомления.
\begin{itemize}
    \item Межпроцедурный статический анализ кода.
          Идея начала активно разрабатываться Томасом Репсом~\sidecite{Reps}.
          Далее последовал ряд, в том числе инженерных работ, применяющих достижимость с контекстно-свободными ограничениями для анализа указателей, анализа алиасов и других прикладных задач~\sidecite{LabelFlowCFLReachability,specificationCFLReachability,Zheng}.
    \item Графовые БД.
          Впервые задача сформулирована Михалисом Яннакакисом~\sidecite{Yannakakis}.
          Запросы с контекстно-свободными ограничениями нашли своё применение различных областях.
          \begin{itemize}
              \item Социальные сети~\sidecite{Hellings2015PathRF}.
              \item RDF обработка~\sidecite{10.1007/978-3-319-46523-4_38}.
              \item Биоинформатика~\sidecite{cfpqBio}.
          \end{itemize}
\end{itemize}

%\begin{itemize}
%    \item OpenCypher~\cite{Kuijpers:2019:ESC:3335783.3335791}
%    \item J.Hellings. CFPQ~\cite{hellingsRelational,hellings2015querying,Hellings2015PathRF}
%    \item Zhang. CFPQ on rdf graphs~\cite{10.1007/978-3-319-46523-4_38}
%    \item Bradford~\cite{bradford2007quickest,ward2008distributed,bradford2016fast,Bradford:2008:LCG:1373936.1373946}
%\end{itemize}


%\section{Вопросы и задачи}
%\begin{enumerate}
%    \item Пусть есть граф. Задайте грамматику для поиска всех путей, таких, что....
%    \item Существует ли в графе !!! путь из А в Б, такой что!!!
%    \item Для графа !!! постройте все пути, удовлетворяющие !!!!
%
%    \item Задача 1
%    \item Задача 2
%\end{enumerate}
